\documentclass[12pt,a4paper,twoside]{article}

\usepackage[utf8]{inputenc}
\usepackage[T1]{fontenc}
\usepackage[catalan]{babel}
\usepackage{amsmath}
\usepackage{amssymb}
\usepackage{geometry}
\usepackage{enumitem}
\usepackage{titlesec}
\usepackage{parskip}
\usepackage{fancyhdr}
\usepackage{tabularx}
\usepackage{booktabs}
\usepackage{array}
\usepackage{tcolorbox}
\usepackage{tikz}

\geometry{top=2.0cm, bottom=1.5cm, left=1.5cm, right=1.5cm, headheight=15pt, headsep=0.5cm}

\pagestyle{fancy}
\fancyhf{}
\fancyhead[L]{\small\textit{Física · 4t ESO · Escola Cervetó}}
\fancyhead[R]{\small\textit{Cinemàtica}}
\fancyfoot[C]{\small\thepage}
\renewcommand{\headrulewidth}{0.4pt}

\titleformat{\section}[block]{\large\bfseries}{}{0em}{}[\vspace{2pt}\hrule\vspace{4pt}]

\newcounter{exercici}
% Exercici command with grade box on the same line
\newcommand{\exercici}[2]{%
  \stepcounter{exercici}
  \section*{Exercici \theexercici\ \textnormal{(#1 punts)} \hfill \normalfont\large Nota: \framebox[2.5cm]{\rule{0pt}{20pt}} / #1}
  \vspace{-10pt}
  \textit{#2}
  \vspace{5pt}
}

% Grade box - removed as it's now integrated in the header
\newcommand{\gradebox}[1]{}

% Plots filling most of the page with more boxes
\newcommand{\drawplots}[1]{%
  \vfill
  \begin{center}
    \begin{tikzpicture}[xscale=2.1, yscale=2.4]
      % Position Grid (Top)
      \draw[step=0.2, gray!50, thin] (0,6.6) grid (8.5,9.1);
      \draw[->, thick] (0,7.85) -- (8.8,7.85) node[right] {$t$};
      \draw[->, thick] (0,6.6) -- (0,9.4) node[above] {$#1$};
      
      % Velocity Grid (Middle)
      \draw[step=0.2, gray!50, thin] (0,3.3) grid (8.5,5.8);
      \draw[->, thick] (0,4.55) -- (8.8,4.55) node[right] {$t$};
      \draw[->, thick] (0,3.3) -- (0,6.1) node[above] {$v$};

      % Acceleration Grid (Bottom)
      \draw[step=0.2, gray!50, thin] (0,0) grid (8.5,2.5);
      \draw[->, thick] (0,1.25) -- (8.8,1.25) node[right] {$t$};
      \draw[->, thick] (0,0) -- (0,2.8) node[above] {$a$};
    \end{tikzpicture}
  \end{center}
}

\begin{document}

% ================================================================
% PÀGINA 1: INSTRUCCIONS I PUNTUACIÓ
% ================================================================
\begin{center}
  \vspace*{-15pt}
  {\Huge\bfseries Col·lecció 2: Cinemàtica}\\[8pt]
  {\Large Física · 4t ESO \quad·\quad Escola Cervetó \quad·\quad Prof.\ Oriol Farrés}\\[12pt]
  \hrule height 1.5pt
  \vspace{5pt}
\end{center}

\noindent
\begin{tabularx}{\textwidth}{|X|}
  \hline
  \rule{0pt}{15pt}\textbf{\large Nom i cognoms:}\\[5pt]
  \hline
\end{tabularx}

\vspace{0.6cm}

\section*{Instruccions de l'Examen}

\begin{itemize}[leftmargin=1.5em, itemsep=4pt, parsep=0pt]
  \item \textbf{Constants:} Usa $g = 9{,}81\ \mathrm{m/s^2}$. Es permet usar $\mathbf{\frac{1}{2} g = 4{,}9}$ per simplificar.
  \item \textbf{Càlculs:} Arrodoneix a \textbf{2 decimals} a cada pas del procés.
  \item \textbf{Unitats:} \textbf{Mostra sempre la conversió al SI} de forma explícita.
  \item \textbf{Estructura Obligatòria:} Cal seguir els 4 passos per a cada exercici: (1)~Fórmules, (2)~Dades/SI, (3)~Esquema del problema, (4)~Resolució de cada apartat.
  \item \textbf{Gràfiques:} Usa els eixos proporcionats. Deuen ser qualitativament correctes. \textbf{És obligatori indicar les mètriques i els punts rellevants als eixos}. Dibuixar només la forma de la gràfica sense valors és fonamentalment invàlid i es considerarà com no realitzat.
\end{itemize}

\vspace{0.3cm}

\section*{Taula de Conversió de Longitud}
\noindent
\begin{tabularx}{\textwidth}{|*{7}{>{\centering\arraybackslash}X|}}
  \hline
  \rule{0pt}{12pt}\textbf{km} & \textbf{hm} & \textbf{dam} & \textbf{m} & \textbf{dm} & \textbf{cm} & \textbf{mm} \\[2pt]
  \hline
  \rule{0pt}{15pt} 0 & 2 & 0 & 0 & 0 & 0 & 0 \\[5pt]
  \hline
  \rule{0pt}{15pt} 0 & 7 & 5 & 0 & 0 & 0 & 0 \\[5pt]
  \hline
  \rule{0pt}{15pt} & & & & & & \\[5pt]
  \hline
  \rule{0pt}{15pt} & & & & & & \\[5pt]
  \hline
  \rule{0pt}{15pt} & & & & & & \\[5pt]
\end{tabularx}

\vfill

\section*{Graella de Puntuació}
\noindent
\begin{tabularx}{\textwidth}{|X r|}
  \hline
  \rule{0pt}{40pt} \textbf{\Huge Nota final:} & \fontsize{35}{42}\selectfont \hspace{4cm} \textbf{/ 10} \\[15pt]
  \hline
\end{tabularx}

\cleardoublepage

% ================================================================
% PÀGINA 2: EXERCICI 1
% ================================================================
\exercici{2,5}{MRU · Sherlock Holmes i la persecució a Londres}

El 1891, Sherlock Holmes rep una pista sobre el parador de Moriarty als molls de Londres. Holmes surt del 221B de Baker Street a peu a una velocitat constant de $\mathbf{9\text{ km/h}}$. L'amagatall es troba a una distància exacta de $\mathbf{120\text{ decàmetres}}$ de casa seva. 

El Dr.\ Watson, que s'estava posant l'abric, s'adona $\mathbf{4\text{ minuts}}$ més tard que Holmes s'ha deixat la seva famosa lupa i decideix agafar un cotxe de cavalls per atrapar-lo. El cotxe de cavalls es mou a una velocitat constant de $\mathbf{200\text{ m/minut}}$.

\begin{enumerate}[label=\textbf{\alph*)}]
  \item Escriu les equacions de posició $x(t)$ per a Holmes i Watson.
  \item Arribarà Holmes a l'amagatall abans de $15\text{ minuts}$? Si és que sí, amb quin temps de marge?
  \item Aconseguirà Watson lliurar-li la lupa a Holmes abans que aquest arribi a l'amagatall? Si és que no, a quina distància (en metres) es trobarà Watson de l'amagatall quan Holmes hi entri?
  \item Dibuixa les 3 gràfiques del moviment ($x-t$, $v-t$, $a-t$) per a tots dos personatges en els mateixos eixos.
\end{enumerate}

\drawplots{x}
\gradebox{2,5}

\newpage

% ================================================================
% PÀGINA 3: EXERCICI 2
% ================================================================
\exercici{2,5}{MRUA · Maglev: El tren del futur}

L'any 2023, durant unes proves de velocitat a la Xina, un prototip de tren de levitació magnètica (Maglev) arrenca des del repòs en una via de proves rectilínia. El tren accelera de forma uniforme fins a assolir la seva velocitat de creuer de $\mathbf{540\text{ km/h}}$ en un temps de $\mathbf{0,5\text{ minuts}}$. 

Just en arribar a aquesta velocitat, el sistema de seguretat detecta un obstacle a la via i el conductor activa el fre magnètic, aturant el tren completament i de manera uniforme en una distància de $\mathbf{25\,000\text{ decímetres}}$.

\begin{enumerate}[label=\textbf{\alph*)}]
  \item Quina ha estat l'acceleració (en $\text{m/s}^2$) en cadascun dels dos trams ($a_1$ i $a_2$)?
  \item Escriu les equacions del moviment per a cadascun dels dos trams.
  \item Quina distància total (en quilòmetres) ha recorregut el tren des de l'inici fins que s'ha aturat definitivament?
  \item Dibuixa les 3 gràfiques ($x-t$, $v-t$, $a-t$) del moviment complet.
\end{enumerate}

\drawplots{x}
\gradebox{2,5}

\newpage

% ================================================================
% PÀGINA 4: EXERCICI 3
% ================================================================
\exercici{2,5}{Caiguda Lliure · Leonardo da Vinci i la Catedral de Milà}

Diu la llegenda que Leonardo da Vinci, mentre observava la construcció de la catedral, volia estudiar el pes de l'aire. Des d'una de les agulles de la catedral, situada a $\mathbf{6,2\text{ decàmetres}}$ d'altura, deixa caure una esfera de bronze ($v_0=0$). 

Exactament en el mateix instant, un aprenent situat en un nivell inferior, a $\mathbf{32\,000\text{ mil·límetres}}$ d'altura, llança una segona esfera verticalment cap avall amb una velocitat inicial de $\mathbf{900\text{ cm/s}}$.

\begin{enumerate}[label=\textbf{\alph*)}]
  \item Escriu les equacions de posició $y(t)$ per a les dues esferes.
  \item Arribaran a xocar les dues esferes a l'aire abans d'impactar contra el terra? Si és que sí, calcula l'altura (en metres) i l'instant exacte del xoc.
  \item Si no xoquen, quin dels dos objectes arribarà primer al terra i amb quina velocitat (en $\text{km/h}$) impactarà cadascun?
  \item Dibuixa les 3 gràfiques ($y-t$, $v-t$, $a-t$) d'ambdós objectes.
\end{enumerate}

\drawplots{y}
\gradebox{2,5}

\newpage

% ================================================================
% PÀGINA 5: EXERCICI 4
% ================================================================
\exercici{2,5}{Tir Vertical · El llançament olímpic de bàsquet}

Durant la final dels Jocs Olímpics, un jugador de bàsquet llança la pilota verticalment cap amunt des de la seva mà, situada a una altura inicial de $\mathbf{215\text{ centímetres}}$ de la pista, per celebrar la victòria. Imprimeix a la pilota una velocitat inicial de $\mathbf{21,6\text{ km/h}}$. La pilota puja fins al seu punt més alt i després cau lliurement fins a tocar el parquet de la pista.

\begin{enumerate}[label=\textbf{\alph*)}]
  \item Escriu l'equació de posició $y(t)$ de la pilota.
  \item A quina altura màxima, respecte al terra de la pista, arribarà la pilota ($y_{\max}$) i en quin instant de temps s'assoleix aquest punt ($t_{\max}$)?
  \item Quant de temps (en segons) estarà la pilota en l'aire des que surt de la mà fins que toca el terra per primer cop?
  \item Amb quina velocitat (en $\text{m/s}$) impacta la pilota contra el parquet?
  \item Dibuixa les 3 gràfiques ($y-t$, $v-t$, $a-t$) de l'acció completa, indicant-hi els punts rellevants.
\end{enumerate}

\drawplots{y}
\gradebox{2,5}

\newpage

% ================================================================
% PÀGINA 6: EXERCICI 5
% ================================================================
\exercici{2,5}{MRU · Marie Curie i el transport de ràdi}

L'any 1914, durant la Primera Guerra Mundial, Marie Curie condueix una unitat mòbil de radiografia (les ``Petites Curies'') cap al front de batalla. Surt de l'Institut del Radi a una velocitat constant de $\mathbf{54\text{ km/h}}$. L'hospital de campanya es troba a una distància de $\mathbf{180\text{ hectòmetres}}$ de París.

Un comboi militar de suport, que ha d'escortar la unitat, s'adona $\mathbf{10\text{ minuts}}$ més tard que la Marie ja ha marxat i surt immediatament darrere seu a una velocitat constant de $\mathbf{1\,200\text{ m/minut}}$.

\begin{enumerate}[label=\textbf{\alph*)}]
  \item Escriu les equacions de posició $x(t)$ per a la Marie Curie i per al comboi d'escorta.
  \item Arribarà la Marie a l'hospital abans de $25\text{ minuts}$? Si és que sí, amb quin temps de marge (en minuts)?
  \item Aconseguirà el comboi atrapar la unitat mòbil abans que arribi a la destinació? Si és que no, a quina distància (en metres) es trobarà el comboi de l'hospital quan la Marie hi entri?
  \item Dibuixa les 3 gràfiques del moviment ($x-t$, $v-t$, $a-t$) per a tots dos vehicles en els mateixos eixos.
\end{enumerate}

\drawplots{x}
\gradebox{2,5}

\newpage

% ================================================================
% PÀGINA 7: EXERCICI 6
% ================================================================
\exercici{2,5}{MRUA · Falcon 9 i el retorn a la Terra}

Durant una missió de SpaceX, el coet Falcon 9 realitza una maniobra de prova en una plataforma horitzontal. El coet arrenca des del repòs i accelera constantment fins a assolir una velocitat de $\mathbf{720\text{ km/h}}$ en un temps de $\mathbf{0,25\text{ minuts}}$.

En arribar a aquesta velocitat, s'activa el protocol de seguretat per un error en els sensors de telemetria. El motor de retropropulsió entra en funcionament, aturant el coet completament i de manera uniforme en un espai de $\mathbf{150\,000\text{ decímetres}}$.

\begin{enumerate}[label=\textbf{\alph*)}]
  \item Quina ha estat l'acceleració (en $\text{m/s}^2$) en cadascun dels dos trams ($a_1$ i $a_2$)?
  \item Escriu les equacions del moviment per a cadascun dels dos trams.
  \item Quina distància total (en quilòmetres) ha recorregut des de l'arrencada fins a la detenció total?
  \item Dibuixa les 3 gràfiques ($x-t$, $v-t$, $a-t$) de tot el moviment complet (unint els dos trams).
\end{enumerate}

\drawplots{x}
\gradebox{2,5}

\newpage

% ================================================================
% PÀGINA 8: EXERCICI 7
% ================================================================
\exercici{2,5}{Caiguda Lliure · L'experiment del Far de Sant Sebastià}

Per verificar les lleis de la gravetat, es realitza un experiment des d'un penya-segat a la costa catalana. Es deixa caure des de la llanterna del far, situada a $\mathbf{9,2\text{ decàmetres}}$ sobre el nivell del mar, una esfera de plom ($v_0 = 0$).

Exactament en el mateix instant, des d'una plataforma d'observació situada més avall, a $\mathbf{55\,000\text{ mil·límetres}}$ d'altura respecte al mar, es llança una segona esfera de fusta verticalment cap avall amb una velocitat inicial de $\mathbf{650\text{ cm/s}}$.

\begin{enumerate}[label=\textbf{\alph*)}]
  \item Escriu les equacions de posició $y(t)$ per a totes dues esferes.
  \item Arribaran a xocar les dues esferes a l'aire abans d'impactar contra l'aigua? Si és que sí, calcula l'altura (en metres) i l'instant de temps de l'impacte.
  \item Si no xoquen, quina esfera arribarà primer a l'aigua i amb quina velocitat (en $\text{km/h}$) ho farà cada una?
  \item Dibuixa les 3 gràfiques ($y-t$, $v-t$, $a-t$) d'ambdós objectes.
\end{enumerate}

\drawplots{y}
\gradebox{2,5}

\newpage

% ================================================================
% PÀGINA 9: EXERCICI 8
% ================================================================
\exercici{2,5}{Tir Vertical · L'erupció de l'Etna}

Durant una petita explosió estromboliana al volcà Etna, un fragment de lava (bomba volcànica) és projectat verticalment cap amunt des d'una fissura situada a una altura de $\mathbf{1\,200\text{ centímetres}}$ respecte a la base del cràter. El fragment surt amb una velocitat inicial de $\mathbf{108\text{ km/h}}$. El fragment puja fins al seu punt màxim, on la velocitat s'anul·la, i després cau lliurement fins a la base del cràter.

\begin{enumerate}[label=\textbf{\alph*)}]
  \item Escriu l'equació de posició $y(t)$ del fragment de lava.
  \item A quina altura màxima, respecte a la base del cràter, arribarà el fragment ($y_{\max}$) i en quin instant s'assoleix aquest punt ($t_{\max}$)?
  \item Quant de temps (en segons) estarà volant el fragment des que surt de la fissura fins que impacta contra la base del cràter?
  \item Amb quina velocitat impacta el fragment contra el terra?
  \item Dibuixa les 3 gràfiques ($y-t$, $v-t$, $a-t$) de l'acció completa, explicant els punts i instants rellevants.
\end{enumerate}

\drawplots{y}
\gradebox{2,5}

\newpage

% ================================================================
% PÀGINA 10: DEMOSTRACIONS I EXTRA
% ================================================================
\section*{Punts Extra: Demostracions Teòriques}

\textbf{Instruccions:} Per a cadascuna de les següents deduccions, primer explica sota quines condicions són aplicables i després desenvolupa el procés matemàtic pas a pas.

\vspace{10pt}

\section*{Exercici Extra 1 \textnormal{(2 punts)} \hfill \normalsize\normalfont\textit{Deducció de fórmules}}

A partir de les equacions generals del MRUA per a un llançament vertical cap amunt:

\begin{enumerate}[label=\textbf{\alph*)}, itemsep=15pt]
  \item \textbf{Temps Màxim ($t_{\max}$):} Demostra que el temps necessari per arribar al punt més alt del llançament es pot calcular com $t_{\max} = \frac{v_0}{g}$. \textit{(1 punt)}
  \item \textbf{Altura Màxima ($y_{\max}$):} Utilitzant el resultat anterior, demostra que l'altura màxima assolida (partint d'una altura inicial $\mathbf{y_0 = h}$) és $y_{\max} = h + \frac{v_0^2}{2g}$. \textit{(1 punt)}
\end{enumerate}

\vfill

\noindent
\begin{tabularx}{\textwidth}{|>{\raggedright\arraybackslash}X|>{\centering\arraybackslash}p{5cm}|}
  \hline
  \rule{0pt}{22pt}\textbf{Demostració A ($t_{\max}$)} & \Large\hspace{2.5cm}~/~1,0 \\[8pt]
  \hline
  \rule{0pt}{22pt}\textbf{Demostració B ($y_{\max}$)} & \Large\hspace{2.5cm}~/~1,0 \\[8pt]
  \hline
  \rule{0pt}{22pt}\textbf{\Large TOTAL EXTRA} & \Large\textbf{\hspace{2.5cm}~/~2,0} \\[8pt]
  \hline
\end{tabularx}

\cleardoublepage

% ================================================================
% PÀGINES EXTRA: ESBORRANY
% ================================================================
\pagestyle{empty}
\begin{center}
  \vspace*{2cm}
  {\Large\bfseries ESPAI PER A CÀLCULS I ESBORRANY}\\[1cm]
  \textit{(Aquestes pàgines no seran avaluades si no s'indica el contrari)}
\end{center}

\newpage
\mbox{}

\end{document}
